\section{Introduction}
\label{sec:intro}

The electrocardiogram (ECG) represents the primary non-invasive diagnostic modality for assessing cardiac electrophysiology and detecting cardiovascular diseases (CVDs), which remain the leading cause of global mortality \cite{clifford2006}. In recent years, continuous cardiac monitoring has undergone a paradigm shift from specialized hospital environments toward ambulatory and wearable consumer electronics, such as smartwatches, Holter patches, and chest straps. However, in mobile and unconstrained settings, acquired ECG signals are intrinsically vulnerable to multiple sources of physiological and environmental contamination \cite{clifford2006,jin2024}. The most prevalent artifacts include baseline wander (BW), induced by respiration and thoracic impedance variations ($0.05$--$1$\,Hz); muscle artifacts (MA), generated by skeletal muscle contraction ($5$--$50$\,Hz); and electrode motion (EM), caused by movement between the skin and sensor interface. Because MA and EM artifacts overlap significantly with the diagnostic frequency bands of the P-wave, QRS complex, and T-wave, simple frequency-domain filtering cannot isolate noise without corrupting essential clinical wave morphologies.

To address non-stationary cardiac noise, traditional techniques such as digital bandpass filters, adaptive filtering, empirical mode decomposition (EMD), and discrete wavelet thresholding have been extensively utilized \cite{clifford2006,goldberger2000}. Although computationally inexpensive, these classical approaches depend heavily on handcrafted heuristic rules, mother wavelet selection, and static thresholding policies. Consequently, aggressive filtering frequently suppresses low-amplitude diagnostic waveforms (such as P-wave onset or subtle ST-segment deviations) or induces non-linear phase shifts that distort the spatial positioning of R-peaks.

\begin{figure*}[!t]
\centering
\begin{adjustbox}{max width=0.98\textwidth}
\begin{tikzpicture}[
  node distance=4mm,
  blk/.style={draw, rectangle, rounded corners=2pt, minimum height=7mm, minimum width=1.4cm, align=center, font=\footnotesize, fill=white, line width=0.6pt},
  den/.style={blk, fill=green!8, draw=green!60!black},
  cls/.style={blk, fill=blue!8, draw=blue!60!black},
  qnt/.style={blk, fill=violet!10, draw=violet!70!black},
  lossn/.style={blk, fill=orange!10, draw=orange!70!black},
  arr/.style={-{Stealth[length=2mm]}, thick},
  darr/.style={-{Stealth[length=2mm]}, thick, dashed},
  grp/.style={draw=black!60, dashed, rounded corners=4pt, inner sep=3.5mm, font=\scriptsize\bfseries}
]

% ==================== STAGE 1 ====================
\node[blk] (n1) {Noisy Segment $x[n]$\\($8192$ samples, $\mathrm{mV}$)};
\node[den, right=6mm of n1] (n2) {\textbf{\HSL}\\\textit{Wavelet Denoiser}};
\node[blk, right=6mm of n2] (n3) {Denoised Output\\$\hat{y}[n] = x + \Delta$};
\node[lossn, right=8mm of n3] (ls) {Mixed Loss $\mathcal{L}_{\mathrm{mixed}}$\\$0.85\mathcal{L}_{\mathrm{MSE}} + 0.15\mathcal{L}_{L_1}$};
\node[blk, below=4mm of ls] (n4) {Clean Target $x_0[n]$\\($8192$ samples, $\mathrm{mV}$)};

\draw[arr] (n1) -- (n2);
\draw[arr] (n2) -- (n3);
\draw[arr] (n3) -- (ls);
\draw[arr] (n4) -- (ls);

\begin{scope}[on background layer]
\node[grp, fit=(n1)(n2)(n3)(ls)(n4), label={[font=\footnotesize\bfseries, anchor=south west]north west:\textbf{Stage 1: Upstream Denoising (\HSL{} Trained Here)}}] (g1) {};
\end{scope}

% ==================== STAGE 2 ====================
\node[blk, below=26mm of n1] (m1) {Segmented Beat $x_d$\\($256$ samples, R-peak)};
\node[den, right=5mm of m1] (m2) {\textbf{\HSL}\\(Weights Frozen)};
\node[cls, right=5mm of m2] (m3) {Adaptive\\Z-Score Norm};
\node[cls, right=5mm of m3] (m4) {1D ResNet\\Morphology (128D)};
\node[cls, right=5mm of m4] (m5) {Feature Fusion\\(160D)};
\node[qnt, right=5mm of m5] (m6) {Classification Head\\(Classical MLP / VQC)};
\node[blk, right=5mm of m6] (m7) {Diagnostic Output\\\textbf{N / S / V}};

% RR Context Branch
\node[cls, below=6mm of m3] (r1) {6D RR-Intervals\\($f_1, \dots, f_6$)};
\node[cls, right=5mm of r1] (r2) {BN-Linear-GELU\\Rhythm Pathway (32D)};

\draw[arr] (m1) -- (m2);
\draw[arr] (m2) -- (m3);
\draw[arr] (m3) -- (m4);
\draw[arr] (m4) -- (m5);
\draw[arr] (m5) -- (m6);
\draw[arr] (m6) -- (m7);

\draw[arr] (r1) -- (r2);
\draw[arr] (r2.east) -| (m5.south);

% Frozen weights arrow
\draw[darr] (n2.south) -- ++(0,-7mm) -| node[pos=0.25, above, font=\scriptsize\bfseries, text=black!80]{Weights Frozen \& Reused} (m2.north);

\begin{scope}[on background layer]
\node[grp, fit=(m1)(m2)(m7)(r1)(r2), label={[font=\footnotesize\bfseries, anchor=south west]north west:\textbf{Stage 2: Downstream Heartbeat Arrhythmia Classification (\RRN{} Trained on DS1/DS2)}}] (g2) {};
\end{scope}

\end{tikzpicture}
\end{adjustbox}
\caption{End-to-end framework for wearable ECG analysis. \textbf{Stage 1 (Top)}: \HSL{} is trained on 8192-sample raw-scale continuous segments under single and multi-source NSTDB noise with a physics-preserving mixed objective. \textbf{Stage 2 (Bottom)}: The frozen denoiser purifies segmented 256-sample cardiac beats, processed by \RRN{} alongside an independent 6-dimensional RR interval timing pathway. Dual interchangeable heads (Classical MLP and 6-qubit Hybrid VQC) perform final 3-class AAMI arrhythmia classification (N, S, V) on the 49{,}298-beat inter-patient DS2 test set.}
\label{fig:pipeline}
\end{figure*}

With the advent of deep learning, neural networks have demonstrated superior non-linear mapping capabilities by learning directly from data pairs \cite{vincent2008,chiang2019}. Denoising autoencoders (DAEs), fully convolutional networks (FCNs), and lightweight U-Nets \cite{ronneberger2015,xiong2016} have shown significant noise suppression. More recently, two-dimensional time-frequency representations, such as the short-time Fourier transform (STFT), have been explored to expose multi-scale spectral patterns. However, STFT-based models exhibit substantial computational redundancy due to overlapping analysis windows, and denoising magnitude spectrograms requires separate, computationally intensive phase-retrieval algorithms (e.g., Griffin-Lim), which impedes real-time deployment on battery-constrained wearable edge devices \cite{jin2024}.

Hybrid architectures combining the discrete wavelet transform (DWT) with deep convolutional backbones offer an appealing compromise \cite{jin2024,nguyen2026liwave,nguyen2026dwse}. By embedding invertible, orthogonal wavelet operations within the network, multiresolution subbands are extracted without phase ambiguity or spatial redundancy. Despite these advances, two fundamental shortcomings persist in the state of the art:
\begin{enumerate}
    \item \textbf{High-Frequency Subband Attenuation and Loss Distortion:} Conventional DWT-based models predominantly decompose only the low-frequency approximation branch ($cA$), leaving the high-frequency detail subbands ($cD$) coarsely resolved. Concurrently, standard networks are optimized using Mean Squared Error (MSE), which penalizes large amplitude deviations quadratically. In ECG processing, this causes catastrophic peak-blunting: the network minimizes overall variance by smoothing the sharp R-peak and blunting the steep slopes of the QRS complex, compromising subsequent clinical analysis.
    \item \textbf{Disconnection from Downstream Arrhythmia Diagnosis:} The vast majority of ECG denoising studies evaluate performance exclusively using generic reconstruction metrics, such as Signal-to-Noise Ratio improvement ($\Delta$SNR) and Percentage Root-Mean-Square Difference (PRD). However, a high mathematical $\Delta$SNR does not ensure clinical fidelity. Crucially, denoising models are rarely subjected to downstream arrhythmia classification under the strict de Chazal inter-patient protocol (DS1/DS2) \cite{dechazal2004}, leaving the clinical utility of deep denoisers unverified on unseen patients.
\end{enumerate}

To bridge this divide, this paper proposes an integrated, two-stage clinical ECG processing framework (Fig.~\ref{fig:pipeline}). In Stage~1, we introduce \textbf{\HSL}, an ultra-compact 75{,}133-parameter 1D wavelet residual network that integrates multi-scale Haar DWT/IDWT operators, unified standard residual blocks with squeeze-and-excitation (\texttt{ResSEBlock1D}) with SiLU activations, and high-frequency detail subband residual filtering, trained under a composite time-domain mixed loss objective. In Stage~2, we design \textbf{\RRN}, a downstream heartbeat classifier that takes beats denoised by the frozen \HSL{} model, extracts deep morphological representations via a 1D ResNet encoder (with SiLU activations), fuses an independent 6-dimensional RR contextual rhythm pathway, and evaluates both a classical MLP head and a 6-qubit parameterized variational quantum circuit (VQC) head.

The primary contributions of this work are summarized as follows:
\begin{itemize}
    \item \textbf{Lightweight Wavelet Denoiser (\HSL):} We engineer a 75.1\,K-parameter symmetric wavelet U-Net combining two-level Haar DWT resolution scaling with unified \texttt{ResSEBlock1D} attention and detail subband filtering. \HSL{} achieves an average $\Delta$SNR of \best{10.01\,dB} and PRD of \best{33.88\%} across comprehensive test conditions on MITDB and NSTDB, surpassing recent SOTA DeepFilter \cite{romero2021} (9.64\,dB) and LiWave \cite{nguyen2026liwave} (9.21\,dB) while running at 35.20\,ms per segment on edge CPU.
    \item \textbf{Fiducial-Preserving Mixed Loss Formulation:} We systematically investigate loss objectives and demonstrate that our time-domain mixed loss ($\alpha L_{\mathrm{MSE}} + (1-\alpha)L_{L_1}$, $\alpha=0.8$) prevents the R-peak blunting inherent to pure MSE training, reducing R-amplitude error to only 0.048\,mV and preserving QRS and QT interval durations under NeuroKit2 clinical evaluation.
    \item \textbf{Downstream Inter-Patient Clinical Validation (\RRN):} We couple the frozen denoiser to a downstream classifier and rigorously benchmark on the 49{,}298-beat inter-patient DS2 test set (22 completely unseen patients). Our 5-seed soft-voting ensemble attains \best{96.43\%} accuracy and \best{0.8407} Macro-F1, significantly elevating supraventricular ectopic beat (S-class) positive predictivity (+P) to \best{72.97\%} (+23.27\% over Mond{\'e}jar-Guerra \emph{et al.} \cite{mondejar2019}).
    \item \textbf{Exploration of Hybrid Quantum Computing and Ablation Insights:} We construct and validate a 6-qubit, 24-parameter VQC classifier with classical residual bypass, achieving an ensemble Macro-F1 of \best{0.8176}. Through controlled ablations, we provide the biomedical insight that channel attention (SE1D) is structurally superior to spatial attention (CBAM1D) for 1D ECG waveforms.
\end{itemize}

The remainder of this article is structured as follows. Section~\ref{sec:related} reviews related work and provides a comprehensive taxonomic comparison. Section~\ref{sec:method} details the mathematical formulation of \HSL{} and \RRN. Section~\ref{sec:exp} describes the dataset, experimental configurations, and metrics. Section~\ref{sec:results} presents quantitative denoising benchmarks, morphology preservation, latency, and classification results. Section~\ref{sec:discussion} discusses biological implications, quantum modeling, and edge deployment trade-offs. Section~\ref{sec:conclusion} concludes the paper.
